\documentclass[a4paper, 12pt]{extreport}


\usepackage{fontenc}
\usepackage{tempora}
\usepackage[utf8]{inputenc}
\usepackage[english,russian]{babel}
\usepackage{amssymb,amsfonts,amsmath,mathtext,cite,enumerate,float}
\usepackage{pgfplots}
\usepackage{graphicx}
\usepackage{tocloft}
\usepackage{listings}
\usepackage[tableposition=top]{caption}
\usepackage{titlesec}
\usepackage{setspace}
\usepackage{geometry}
\usepackage{indentfirst}
\usepackage{fancyhdr}
\usepackage{etoolbox}
\usepackage{pdfpages}
\usepackage{enumerate,letltxmacro}
\usepackage{threeparttable}
\usepackage[hidelinks]{hyperref}
\usepackage{flafter}
\usepackage{enumitem}
\usepackage{multirow}
\usepackage[figure,table]{totalcount}
\usepackage{lastpage}
\usepackage{longtable}
\usepackage{array}
\usepackage[most]{tcolorbox}
\usepackage{xcolor}
\usepackage{xparse}
\usepackage{textcomp}


\setlist{nosep}

\setcounter{secnumdepth}{3}


\newcommand{\ssr}[1]{
	\begin{center}
		\LARGE\bfseries{#1}
	\end{center} 
	\addcontentsline{toc}{chapter}{#1}  
}

\newcommand{\ssrn}[1]{%
  \begin{center}
    \LARGE\bfseries #1
  \end{center}%
   \vspace{1em}
}





\DeclareUnicodeCharacter{2713}{}  % ✓
\DeclareUnicodeCharacter{2717}{X} % ✗
\DeclareUnicodeCharacter{274C}{X} % ❌
\DeclareUnicodeCharacter{1F4E6}{} % 📦
\DeclareUnicodeCharacter{2705}{}  % ✅
\DeclareUnicodeCharacter{1F50D}{} % 🔍
\DeclareUnicodeCharacter{1F4CC}{} % 📌
\DeclareUnicodeCharacter{200A}{~}
\DeclareUnicodeCharacter{202F}{~}





\newcommand{\russianAppendixLetter}[1]{%
  \ifcase#1%
  \or А\or Б\or В\or Г\or Д\or Е\or Ж\or И\or К\or Л\or М\or Н\or П\or Р\or С\or Т\or У\or Ф\or Х\or Ц\or Ш\or Щ\or Э\or Ю\or Я%
  \else\@ctrerr % Ошибка, если счётчик вышел за пределы доступных букв
  \fi
}

% Команда для начала раздела приложений
\newcommand{\BeginAppendices}{%
  \clearpage % Начать приложения с новой страницы
  \appendix % Стандартная команда LaTeX для переключения в режим приложений
            % Она сбрасывает счётчик глав и меняет \thechapter на \Alph{chapter}

  % Переопределяем отображение номера главы на нашу русскую букву
  % \appendix уже сбросил chapter в 0. \refstepcounter{chapter} в \prilojenie сделает его 1 для первого приложения.
  \renewcommand{\thechapter}{\russianAppendixLetter{\value{chapter}}}

  % Меняем название "Глава" на "Приложение" для заголовков и TOC
  % Для классов report и book это \chaptername. Но \appendix может уже менять \chaptername на \appendixname.
  % \renewcommand{\chaptername}{Приложение} % Это изменит "Глава X" на "Приложение X"
  % Однако, мы хотим "Приложение А", "Приложение Б" в заголовке, а не "Приложение А Название"
  % Поэтому мы будем форматировать заголовок полностью в команде \prilojenie.
  % Но для TOC это может быть полезно, если стиль TOC использует \chaptername.
  % Вместо этого, \addcontentsline будет формировать нужную строку.

  % Переопределяем нумерацию для секций, рисунков, таблиц, листингов, чтобы они включали букву приложения
  \renewcommand{\thesection}{\thechapter.\arabic{section}}
  \renewcommand{\thesubsection}{\thesection.\arabic{subsection}}
  \renewcommand{\thefigure}{\thechapter.\arabic{figure}}
  \renewcommand{\thetable}{\thechapter.\arabic{table}}
  \renewcommand{\thelstlisting}{\thechapter.\arabic{lstlisting}} % КЛЮЧЕВОЕ для листингов

  % Убедимся, что имя для листинга остаётся "Листинг" (буква А будет добавлена через \thelstlisting)
  \renewcommand{\lstlistingname}{Листинг}

  % Можно добавить общий заголовок "Приложения" в TOC (необязательно)
  % \addcontentsline{toc}{part}{Приложения} % Как "Часть"
  % или использовать \part*{Приложения} для заголовка на отдельной странице
}

% Ваша команда \prilojenie, адаптированная
% Название и необязательный статус приложения.
\newcommand{\prilojenie}[2][справочное]{%
  % \chapter{#1} % Стандартная команда \chapter теперь будет работать как надо,
                % но она выведет "Приложение А: Название приложения" или похоже,
                % в зависимости от стиля. Нам нужно customnoe форматирование.

  \clearpage % Каждое приложение с новой страницы (если \chapter этого не делает в вашем стиле)
  \refstepcounter{chapter} % Увеличиваем счётчик глав (теперь это наш счётчик приложений А, Б, В...)
                           % \refstepcounter также делает метку для \label/\ref
                           % Важно: \appendix сбрасывает chapter в 0, так что первое приложение будет "1" -> "А"

  % Форматирование заголовка приложения
  \phantomsection % Для корректной работы ссылок hyperref
  % \addtocontents{toc}{\protect\addvspace{10pt}} % Добавить немного места в TOC перед записями приложений (если нужно)
  
  \begin{center}
    \textbf{ПРИЛОЖЕНИЕ \thechapter}\\
    (#1)\\
    \textbf{#2}
  \end{center}
  
  % Добавляем в содержание (Table of Contents)
  % Формат: "Приложение А. Название приложения"
  \addcontentsline{toc}{chapter}{Приложение \thechapter\space (#1) #2}
  
  % Установка заголовков для страниц (headers/footers)
  \markboth{\MakeUppercase{Приложение \thechapter\space #2}}{\MakeUppercase{Приложение \thechapter\space #2}}

  % После \refstepcounter{chapter}, счётчики, зависимые от главы (секции, рисунки, таблицы, листинги),
  % должны автоматически сброситься в 0, если они так настроены (обычно так и есть).
  % Например, счётчик lstlisting сбрасывается при каждом \chapter.
}

\newcommand{\abzend}{\newline\indent}

% Структурный элемент отчёта: заголовок без номера, с записью в содержание.
\newcommand{\structel}[1]{%
  \clearpage\phantomsection\addcontentsline{toc}{chapter}{#1}%
  \chapter*{#1}}

\makeatletter
\renewcommand\LARGE{\@setfontsize\LARGE{22pt}{20}}
\renewcommand\Large{\@setfontsize\Large{20pt}{20}}
\renewcommand\large{\@setfontsize\large{16pt}{20}}
\makeatother
\RequirePackage{titlesec}
\newcommand{\gostheading}{\normalfont\normalsize\hyphenpenalty=10000\exhyphenpenalty=10000\relax}
\titleformat{\chapter}[block]{\gostheading\bfseries\filright}{\thechapter}{0.5em}{}
\titleformat{name=\chapter,numberless}[block]{\gostheading\bfseries\filcenter}{}{0pt}{}
\titleformat{\section}[block]{\gostheading\bfseries\filright}{\thesection}{0.5em}{}
\titleformat{\subsection}[block]{\gostheading\bfseries\filright}{\thesubsection}{0.5em}{}
\titleformat{\subsubsection}[block]{\gostheading\bfseries\filright}{\thesubsubsection}{0.5em}{}
\titlespacing{\chapter}{\parindent}{0pt}{\baselineskip}
\titlespacing{name=\chapter,numberless}{0pt}{0pt}{\baselineskip}
\titlespacing{\section}{\parindent}{\baselineskip}{\baselineskip}
\titlespacing{\subsection}{\parindent}{\baselineskip}{\baselineskip}
\titlespacing{\subsubsection}{\parindent}{\baselineskip}{\baselineskip}
\titlespacing{\paragraph}{\parindent}{10pt}{10pt}

\makeatletter
\renewcommand{\@biblabel}[1]{#1.}
\patchcmd{\thebibliography}{\advance\leftmargin\labelsep}
  {\leftmargin=0pt\itemindent=\dimexpr\parindent+\labelwidth+\labelsep\relax}
  {}{\PackageWarning{preambula}{Bibliography indentation could not be configured}}
\makeatother
%
%\titleformat{\chapter}[hang]{\LARGE\bfseries}{\hspace{1.25cm}\thechapter}{1ex}{\LARGE\bfseries}
%\titleformat{\section}[hang]{\Large\bfseries}{\hspace{1.25cm}\thesection}{1ex}{\Large\bfseries}
%\titleformat{name=\section,numberless}[hang]{\Large\bfseries}{\hspace{1.25cm}}{0pt}{\Large\bfseries}
%\titleformat{\subsection}[hang]{\large\bfseries}{\hspace{1.25cm}\thesubsection}{1ex}{\large\bfseries}
%\titlespacing{\chapter}{0pt}{-\baselineskip}{\baselineskip}
%\titlespacing*{\section}{0pt}{\baselineskip}{\baselineskip}
%\titlespacing*{\subsection}{0pt}{\baselineskip}{\baselineskip}

\geometry{left=30mm}
\geometry{right=15mm}
\geometry{top=20mm}
\geometry{bottom=20mm}
\fancypagestyle{plain}{%
  \fancyhf{}%
  \fancyfootoffset[L]{15mm}%
  \fancyfoot[C]{\thepage}%
  \renewcommand{\headrulewidth}{0pt}%
}
\pagestyle{plain}

\onehalfspacing
\AtBeginDocument{%
  \setlength{\abovedisplayskip}{\baselineskip}%
  \setlength{\belowdisplayskip}{\baselineskip}%
  \setlength{\abovedisplayshortskip}{\baselineskip}%
  \setlength{\belowdisplayshortskip}{\baselineskip}%
}

\renewcommand{\theenumi}{\arabic{enumi}}
\renewcommand{\labelenumi}{\arabic{enumi}\text{)}}
\renewcommand{\theenumii}{.\arabic{enumii}}
\renewcommand{\labelenumii}{\arabic{enumii}\text{)}}
\renewcommand{\theenumiii}{.\arabic{enumiii}}
\renewcommand{\labelenumiii}{\arabic{enumi}.\arabic{enumii}.\arabic{enumiii}.}

\renewcommand{\cftchapleader}{\cftdotfill{\cftdotsep}}
\renewcommand{\cftchapfont}{\normalfont}
\renewcommand{\cftchappagefont}{\normalfont}
\cftsetindents{section}{1em}{2.4em}
\cftsetindents{subsection}{2em}{3.2em}
\cftsetindents{subsubsection}{3em}{4em}
\addto\captionsrussian{%
  \renewcommand{\contentsname}{СОДЕРЖАНИЕ}%
  \renewcommand{\bibname}{СПИСОК ИСПОЛЬЗОВАННЫХ ИСТОЧНИКОВ}%
}

%\addto\captionsrussian{\renewcommand{\figurename}{Рисунок}}
%\DeclareCaptionLabelSeparator{dash}{~---~}
%\captionsetup{labelsep=dash}

%\captionsetup[figure]{justification=centering,labelsep=dash}
%\captionsetup[table]{labelsep=dash,justification=raggedright,singlelinecheck=off}


% НАСТРОЙКА РИСУНКОВ И ТАБЛИЦ

\usepackage{subcaption}
\DeclareCaptionLabelFormat{gostfigure}{Рисунок #2}
\DeclareCaptionLabelFormat{gosttable}{Таблица #2}
\DeclareCaptionLabelSeparator{gost}{~---~}
\captionsetup{labelsep=gost,font={stretch=1}}
\captionsetup[figure]{labelformat=gostfigure,justification=centering}
\captionsetup[table]{labelformat=gosttable,justification=raggedright,singlelinecheck=off}
\renewcommand{\thesubfigure}{\asbuk{subfigure}}




\graphicspath{{images/}{../code/res}}%путь к рисункам

\newcommand{\floor}[1]{\lfloor #1 \rfloor}


\lstdefinestyle{basestyle}{ %
	basicstyle=\ttfamily\normalsize, % размер и начертание шрифта для подсветки кода
	upquote=true,
	numbers=left,               % где поставить нумерацию строк (слева\справа)
	numberstyle=\tiny,           % размер шрифта для номеров строк
	stepnumber=1,                   % размер шага между двумя номерами строк
	numbersep=7pt,                % как далеко отстоят номера строк от подсвечиваемого кода
	xleftmargin=15pt,        % Отступ слева, чтобы номера строк были внутри рамки
    framexleftmargin=15pt,    % Расширение рамки слева, чтобы номера строк были внутри
	showspaces=false,            % показывать или нет пробелы специальными отступами
	showstringspaces=false,      % показывать или нет пробелы в строках
	showtabs=false,             % показывать или нет табуляцию в строках
%	frame=single,              % рисовать рамку вокруг кода
	tabsize=4,                 % размер табуляции по умолчанию равен 2 пробелам
	captionpos=t,              % позиция заголовка вверху [t] или внизу [b] 
	breaklines=true,           % автоматически переносить строки (да\нет)
	breakatwhitespace=false, % переносить строки только если есть пробел
	escapeinside={\#*}{*)},   % если нужно добавить комментарии в коде
	abovecaptionskip=-5pt,
	inputencoding=utf8, 
	extendedchars=\true,      % вроде было нужно для отображения ru символов
	keepspaces=true,           % вроде было нужно для отображения ru символов
	literate={=}{{\char61}}1 {×}{{$\times$}}1 {≠}{{$\ne$}}1 {─}{{\rule[0.5ex]{0.6em}{0.35pt}}}1
             {«}{{\guillemotleft}}1            % Кавычка-ёлочка левая
             {»}{{\guillemotright}}1           % Кавычка-ёлочка правая
             {…}{{\dots}}1                     % Многоточие (как один символ Unicode)
             {—}{{\textemdash}}1               % Длинное тире (Em-dash)
             {–}{{\textendash}}1               % Короткое тире (En-dash)
             {→}{{$\rightarrow$}}1             % Стрелка вправо
             {←}{{$\leftarrow$}}1              % Стрелка влево
             {а}{{\char224}}1 {б}{{\char225}}1 {в}{{\char226}}1 {г}{{\char227}}1 {д}{{\char228}}1 {е}{{\char229}}1 {ё}{{\cyryo}}1 {ж}{{\char230}}1 {з}{{\char231}}1 {и}{{\char232}}1 {й}{{\char233}}1 {к}{{\char234}}1 {л}{{\char235}}1 {м}{{\char236}}1 {н}{{\char237}}1 {о}{{\char238}}1 {п}{{\char239}}1 {р}{{\char240}}1 {с}{{\char241}}1 {т}{{\char242}}1 {у}{{\char243}}1 {ф}{{\char244}}1 {х}{{\char245}}1 {ц}{{\char246}}1 {ч}{{\char247}}1 {ш}{{\char248}}1 {щ}{{\char249}}1 {ъ}{{\char250}}1 {ы}{{\char251}}1 {ь}{{\char252}}1 {э}{{\char253}}1 {ю}{{\char254}}1 {я}{{\char255}}1 {А}{{\char192}}1 {Б}{{\char193}}1 {В}{{\char194}}1 {Г}{{\char195}}1 {Д}{{\char196}}1 {Е}{{\char197}}1 {Ё}{{\CYRYO}}1 {Ж}{{\char198}}1 {З}{{\char199}}1 {И}{{\char200}}1 {Й}{{\char201}}1 {К}{{\char202}}1 {Л}{{\char203}}1 {М}{{\char204}}1 {Н}{{\char205}}1 {О}{{\char206}}1 {П}{{\char207}}1 {Р}{{\char208}}1 {С}{{\char209}}1 {Т}{{\char210}}1 {У}{{\char211}}1 {Ф}{{\char212}}1 {Х}{{\char213}}1 {Ц}{{\char214}}1 {Ч}{{\char215}}1 {Ш}{{\char216}}1 {Щ}{{\char217}}1 {Ъ}{{\char218}}1 {Ы}{{\char219}}1 {Ь}{{\char220}}1 {Э}{{\char221}}1 {Ю}{{\char222}}1 {Я}{{\char223}}1
}

\lstdefinestyle{pythonstyle}{
    style=basestyle,
    language=Python
}

\lstdefinestyle{logstyle}{
    style=basestyle,
    language={},                % Убираем язык, чтобы не подсвечивало синтаксис
    columns=fullflexible,       % Лучшее выравнивание для текстовых логов
    commentstyle=\color{black}, % На всякий случай сбрасываем цвета
    stringstyle=\color{black},
    keywordstyle=\color{black}
}

\lstset{style=basestyle}

% --- 2. Логика на LaTeX3 ---
\ExplSyntaxOn

% Переменные
\tl_new:N \l_mycode_caption_tl
\tl_new:N \l_mycode_label_tl
\tl_new:N \l_mycode_params_tl

\keys_define:nn { mycode } {
    caption .tl_set:N = \l_mycode_caption_tl,
    label .tl_set:N = \l_mycode_label_tl,
    
    unknown .code:n = {
        \tl_if_empty:nTF { #1 }
        { \tl_put_right:Nx \l_mycode_params_tl { , \l_keys_key_tl } }
        { 
            \tl_put_right:Nx \l_mycode_params_tl { , \l_keys_key_tl = }
            \tl_put_right:Nn \l_mycode_params_tl { #1 }
        }
    }
}


% --- 3. Команда \inputcode ---
\NewDocumentCommand{\inputcode}{ +O{} m }{
    % Очистка
    \tl_clear:N \l_mycode_caption_tl
    \tl_clear:N \l_mycode_label_tl
    \tl_clear:N \l_mycode_params_tl
    \tl_set:Nn \l_mycode_params_tl { style=logstyle, aboveskip=0pt, belowskip=0pt }
    
    % Парсинг
    \keys_set:nn { mycode } { #1 }
    
    % Увеличиваем счётчик и ставим label
    \refstepcounter{lstlisting}
    \tl_if_empty:NF \l_mycode_label_tl { \label{\l_mycode_label_tl} }
    
    % --- ГЕНЕРАЦИЯ ЗАГОЛОВКОВ (В ОДНОМ МЕСТЕ) ---
    % 1. Формируем основной заголовок: "Листинг X.Y --- Название"
    \tl_set:Nn \l_mycode_title_main_tl {
        Листинг~\thelstlisting
        \tl_if_empty:NTF \l_mycode_caption_tl {} {~---~\l_mycode_caption_tl }
    }
    
    % 2. Формируем заголовок продолжения"
    \tl_set:Nn \l_mycode_title_cont_tl {
        Листинг~\thelstlisting
        \tl_if_empty:NTF \l_mycode_caption_tl {} {~(продолжение)~---~\l_mycode_caption_tl }
    }
    
    % Рисуем блок
    \tcbinputlisting{
        listing~file={#2},
        listing~only,
        listing~options/.expand~once = \l_mycode_params_tl,
        enhanced, breakable, sharp~corners,
        colback=white, colframe=black, boxrule=0.5pt,
        title={\l_mycode_title_main_tl},
        title~after~break={\l_mycode_title_cont_tl},
        coltitle=black, colbacktitle=white, titlerule=0.5pt,
        fonttitle=\normalfont,
        boxsep=0pt, left=0pt, right=0pt, top=0pt, bottom=0pt,
        lefttitle=3pt, righttitle=3pt, toptitle=2pt, bottomtitle=2pt,
        before~skip=12pt
    }
}
\ExplSyntaxOff

\captionsetup[lstlisting]{singlelinecheck=false, justification=raggedright}

\pgfplotsset{width=0.85\linewidth, height=0.5\columnwidth}

\parindent=1.25cm

%\LetLtxMacro\itemold\item
%\renewcommand{\item}{\itemindent0.75cm\itemold}
% label=\textbullet
\def\labelitemi{---}
\setlist[itemize]{leftmargin=0.65cm, itemindent=0.65cm}
\setlist[enumerate]{leftmargin=0.65cm, itemindent=0.55cm}

\newcommand{\specialcell}[2][c]{%
	\begin{tabular}[#1]{@{}c@{}}#2\end{tabular}}

\newcommand{\n}{\newline}

\frenchspacing

\counterwithout{equation}{chapter}
\counterwithout{table}{chapter}
\begin{document}
\begin{titlepage}
\singlespacing

\noindent\begin{minipage}{0.15\textwidth}
	\includegraphics[width=\linewidth]{b_logo}
\end{minipage}%
\begin{minipage}{0.85\textwidth}\centering
	\textbf{Министерство науки и высшего образования Российской Федерации}\\
	\textbf{Федеральное государственное автономное образовательное учреждение высшего образования}\\
	\textbf{<<Московский государственный технический университет имени Н.Э.~Баумана}\\
	\textbf{(национальный исследовательский университет)>>}\\
	\textbf{(МГТУ им. Н.Э.~Баумана)}
\end{minipage}

\noindent\rule{\linewidth}{3pt}
\newline\newline
\noindent ФАКУЛЬТЕТ \underline{\text{<<Информатика и системы управления>>}} \newline\newline
\noindent КАФЕДРА \underline{\text{<<Программное обеспечение ЭВМ и информационные технологии, ИУ7>>}}

\vspace{1cm}

\begin{center}
	\noindent\begin{minipage}{\textwidth}\centering
		\Large\textbf{Лабораторная работа № 1}\newline
		\textbf{по дисциплине}\newline
		\textbf{<<Программирование (специальные главы)>>}\newline\newline
	\end{minipage}
\end{center}

\noindent\textbf{Тема} Сравнительный анализ решений, сгенерированных большими языковыми моделями\newline\newline
\noindent\textbf{Студент} \underline{\text{Логинов Виктор Николаевич}}\newline\newline
\noindent\textbf{Группа} \underline{\text{СГН3-11М}}\newline\newline
\noindent\textbf{Преподаватели} \underline{\text{Строганов Юрий Владимирович}}\newline

\vfill
\begin{center}
	Москва~\the\year
\end{center}
\end{titlepage}
\setcounter{page}{2}



\structel{РЕФЕРАТ}

Отчет \pageref{LastPage}~с., 1~табл., 9~источн., 7~прил.

{\noindent\raggedright\hyphenpenalty=10000\exhyphenpenalty=10000\relax
БОЛЬШИЕ ЯЗЫКОВЫЕ МОДЕЛИ, ВЕКТОРНЫЕ ПРЕДСТАВЛЕНИЯ, СИМВОЛЬНЫЕ N-ГРАММЫ, СИНТАКСИЧЕСКИЙ АНАЛИЗ, TF--IDF, СРАВНИТЕЛЬНЫЙ АНАЛИЗ\par}

Объектом исследования являются сохранённые ответы пяти больших языковых моделей --- GigaChat, Qwen, Яндекс.Алисы, ассистента МГТУ им.~Н.Э.~Баумана и DeepSeek, содержащие программы обработки русскоязычных документов.

Цель работы --- сравнить, как модели реализуют построение векторных представлений слов и документов на основе символьных n-грамм и синтаксического разбора, и оценить объём исправлений, необходимый для получения проверяемого результата.

Методами исследования служат единый исходный запрос, сохранение ответов без изменений, автоматическая приёмка на семи сценариях, включая полный корпус из 50 документов, и сопоставление числа версий и времени выполнения.

В результате работы финальные версии всех пяти моделей прошли все семь сценариев; число программных версий составило 4 у GigaChat, 3 у Qwen, 2 у Яндекс.Алисы, 3 у ассистента МГТУ и 3 у DeepSeek, а время обработки полного корпуса --- от 5,678 до 10,021~с. Ассистент МГТУ показал наиболее готовую первую реализацию вычислений, Яндекс.Алиса --- наименьшее число исправлений после получения кода, DeepSeek --- наиболее локализуемый последний дефект.

Область применения --- организация цикла разработки программ обработки текста с помощью больших языковых моделей при независимом автоматическом контроле результата.

\clearpage

\renewcommand{\cfttoctitlefont}{\hfill\normalfont\normalsize\bfseries}
\renewcommand{\cftaftertoctitle}{\hfill}
\setlength{\cftbeforetoctitleskip}{0pt}
\setlength{\cftaftertoctitleskip}{\baselineskip}
\clearpage
\tableofcontents

\chapter*{ПЕРЕЧЕНЬ СОКРАЩЕНИЙ И ОБОЗНАЧЕНИЙ}
\addcontentsline{toc}{chapter}{ПЕРЕЧЕНЬ СОКРАЩЕНИЙ И ОБОЗНАЧЕНИЙ}

В настоящем отчёте о лабораторной работе применяют следующие сокращения и обозначения:

\noindent
\begin{tabular}{@{}l@{\quad---\quad}p{0.8\textwidth}@{}}
БЯМ & большая языковая модель \\
CP1251 & кодовая страница Windows-1251 \\
CSR & формат хранения разреженных матриц со сжатием строк \\
DOCX & формат текстового документа Office Open XML \\
HTML & язык гипертекстовой разметки \\
IN & признак синтаксической связи, отнесённый к зависимому токену \\
JSON & текстовый формат обмена структурированными данными \\
NPZ & формат архива массивов NumPy \\
ODT & формат текстового документа OpenDocument \\
OUT & признак синтаксической связи, отнесённый к голове связи \\
PASS & статус успешного прохождения проверки \\
POSIX & интерфейс переносимых операционных систем \\
TF--IDF & вес на основе частоты термина и обратной частоты документа \\
TXT & файл с простым текстом \\
UTF-8 & кодировка символов Unicode
\end{tabular}

\chapter*{ВВЕДЕНИЕ}
\addcontentsline{toc}{chapter}{ВВЕДЕНИЕ}
Работа посвящена получению программы обработки русскоязычных документов посредством диалога с пятью БЯМ. Исследуется весь путь от исходного задания до проверяемого программного результата: генерация, запуск, выявление дефекта и корректирующий запрос. Предмет сравнения — конкретные сохранённые ответы GigaChat, Qwen, Яндекс.Алисы, ассистента МГТУ и DeepSeek, а не универсальный рейтинг этих систем.

Цель эксперимента — установить, какие требования модели реализуют сразу, какие требуют уточнения и насколько полученный результат соответствует независимым проверкам. Для этого используются версии программ, тексты запросов и машинные протоколы. Работой с ответами моделей и технической приёмкой всех версий управлял Codex. Его роль описана в подразделе~\ref{sec:codex}.

\chapter{Аналитическая часть}
\section{Описание задачи и выбранных инструментов}
Исходный корпус включает 50 документов в четырёх форматах: TXT, HTML, DOCX и ODT. Тематические каталоги используются для организации файлов, но не являются входными признаками алгоритма. Программа должна извлечь текст и построить четыре представления: символьные n-граммы слов, TF--IDF документов, синтаксические контексты слов и синтаксические признаки документов.

Эти представления решают разные задачи. Символьные признаки описывают состав слова, TF--IDF отражает распределение лексики между документами, а зависимости учитывают отношения между токенами. Поэтому правильность одной матрицы не доказывает правильность остальных. Кроме чисел, результат обязан содержать подписи строк и столбцов: без них невозможно установить, какому документу, слову или признаку соответствует значение.

Python выбран как общая среда реализации. Для хранения и загрузки разреженных матриц используется SciPy~\cite{scipy}. Расчёт TF--IDF сверяется с документацией scikit-learn~\cite{sklearn}. Синтаксические зависимости определяются средствами spaCy~\cite{spacy}. Проверка чтения документов выделена отдельно: ошибка извлечения текста распространяется на все последующие представления, даже если математические процедуры реализованы верно.

\section{Характеристика используемых БЯМ}
Сервисы рассматриваются как источники программных ответов. Название интерфейса само по себе не задаёт точный идентификатор весов, размер модели или режим декодирования. Поэтому такие характеристики не используются для объяснения результатов, если они не зафиксированы в материалах эксперимента. Практическая характеристика строится по доступному программному результату и истории его исправления.

\subsection{GigaChat}
Описание сервиса GigaChat приведено в документации Сбера~\cite{gigachat}. В архиве представлены четыре версии программы. Наличие полного ответа на каждой итерации позволило проверять не только исправленный фрагмент, но и согласованность импортов, функций чтения и вычислений в целом. Вместе с тем полнота текста не обеспечила готовности к запуску: версия 02 не прошла компиляцию, а версия 03 завершалась ошибкой ещё до выполнения полезной обработки.

Для сравнения это существенный профиль дефектов. Основным препятствием на промежуточных шагах была исполнимость поставленного решения, поэтому оценивать его алгоритмы только по описанию было бы преждевременно. Финальная версия 04 прошла весь набор проверок, но потребовала на один сохранённый ответ больше, чем Qwen, Королёв и DeepSeek.

\subsection{Qwen}
Описание семейства моделей Qwen доступно в официальном репозитории~\cite{qwen}. Для этого сервиса сохранены три программные версии. Две первые проходили сценарий сообщения о повреждённом файле, однако не удовлетворяли остальным сценариям проверки. Таким образом, наличие обработки исключений не означало корректной реализации векторизации и пустых входов.

Особенность этой цепочки — необходимость согласовать подготовку текста с синтаксическим анализом. Удаление пунктуации или иное преобразование текста до spaCy способно изменить дерево зависимостей; совпадение словаря после такой обработки не гарантирует совпадения синтаксических признаков. Финальный результат оценивался уже по всем четырём матрицам, а не по отдельной удачной функции.

\subsection{Яндекс.Алиса}
Яндекс предоставляет доступ к чат-интерфейсу Алиса AI~\cite{alice}. Первая зарегистрированная попытка Алисы завершилась пустым ответом. Это факт отсутствия программного результата, а не синтаксическая ошибка программы. Вторая попытка содержит первую доступную реализацию; третья — её исправление. Поэтому три номера в архиве соответствуют двум версиям кода.

Первая реализация уже проходила оба пустых сценария, но теряла слова без зависимостей и один документ в иной кодировке. Такой результат показывает частичное выполнение требований: корректное поведение на пустом входе сочеталось с неполнотой результата на содержательном корпусе. Для оценки Алисы необходимо различать доступность ответа и качество полученного кода.

\subsection{Ассистент МГТУ им. Н.Э. Баумана}
ИИ-ассистент входит в состав университетской платформы «Королёв ИИ»~\cite{korolev}. Ответы сервиса assistant.korolev.bmstu.ru предоставлены пользователем и сохранены как три версии Королёва. Уже первая версия проходила базовый сценарий и сценарий повторяющихся слов. Это наиболее сильный старт по этим двум содержательным проверкам среди первых доступных реализаций в данном архиве.

Оставшиеся ошибки относились к пустым входам, ODT, кодировке и контракту результата. Вторая версия исправила вычисления на всех пяти контрольных наборах с эталонными матрицами, но ещё не прошла повреждённый файл и полный корпус. Таким образом, последующее уточнение касалось уже не только алгоритма, но и переносимости путей и явности сообщения об ошибке.

\subsection{DeepSeek}
Сервис DeepSeek доступен через веб-чат~\cite{deepseek}. В работе получены три версии программы. Первая формировала NPZ через общий механизм NumPy; такой архив не распознавался штатным загрузчиком разреженных матриц SciPy. Даже наличие массивов значений и индексов внутри файла не выполняет требование совместимости формата.

Вторая версия проходила шесть из семи сценариев, включая полный корпус, но падала на специально подготовленном ODT. Третья исправила импорт элементов таблиц. Эта история особенно полезна для методики: положительный результат на большом корпусе не заменяет целевой тест конструкции, которая в корпусе не проявилась.

\section{Сравнительный анализ выбранных БЯМ}
Сравнение разделяет три характеристики: готовность первого ответа, количество последующих исправлений и итоговое соответствие проверкам. По итоговому статусу различий нет: последние версии всех пяти систем получили PASS. По начальному состоянию различия значительны: Королёв уже решал базовые вычислительные сценарии, Алиса не выдала программу при первой попытке, а у DeepSeek проверку блокировал формат сохранения матриц.

Количество версий характеризует конкретную историю взаимодействия, но не измеряет скорость модели. Для вывода о скорости генерации потребовались бы времена ответов и одинаковые условия сервиса; таких измерений здесь нет. Аналогично нельзя считать большее число исправлений доказательством худшего качества модели на любых задачах: исследована одна постановка и по одной последовательной цепочке ответов.

\chapter{Конструкторский раздел}
\section{Организация работы и приёмки под управлением Codex}
\label{sec:codex}
Внутренним контроллером экспериментального процесса служил программный агент Codex от OpenAI~\cite{codex}. На стороне исследователя он организовывал весь цикл работы с ответами сравниваемых БЯМ и выполнял техническую приёмку каждой версии. Исходный запрос задавал входные форматы, четыре матрицы и правила сохранения. Codex сопоставлял полученный результат с этими требованиями и определял, можно ли перейти к следующему этапу проверки либо требуется новый исправляющий запрос.

После получения ответа Codex сохранял его в каталоге соответствующей итерации и извлекал программный код для запуска. Для каждой версии связывались исходный ответ, извлечённые файлы, команда запуска, журнал выполнения и протокол проверки. Такая фиксация позволяла установить, какая именно программа была проверена и к какому ответу относятся найденные дефекты. Перед испытаниями Codex проверял синтаксис Python, затем запускал программу отдельным процессом на контрольных наборах и исходном корпусе.

Приёмка охватывала весь контракт результата. Codex контролировал состав выходных файлов, упорядоченность меток, размеры и формат матриц, конечность значений и совпадение с эталоном на наборах, для которых он был вычислен. Отдельно проверялись пустые входы, смешанные форматы, диагностика повреждённого документа и сохранение всех 50 документов полного корпуса. Успешный запуск становился одним из оснований решения о версии вместе с результатами содержательных проверок выходных данных.

По результатам проверки Codex локализовал расхождения, фиксировал причину отклонения и составлял корректирующий запрос. В нём указывались наблюдаемый дефект, ожидаемое поведение и критерий следующей проверки. Например, у первой версии DeepSeek штатный загрузчик SciPy не распознал архив разреженной матрицы. В следующем запросе было уточнено требование совместимости с \texttt{scipy.sparse.load\_npz}. Повторная проверка показала, что после устранения этой ошибки отдельно оставалась проблема импорта элементов таблиц ODT. Последующий запрос был посвящён этому дефекту.

Ответы и новые версии программ формировали сравниваемые модели. Codex проверял полученное исправление и повторно выполнял весь применимый набор сценариев, включая ранее успешно пройденные случаи. Он принимал техническое решение о допуске версии к итоговому сравнению, связывая это решение с сохранёнными протоколами. Весь цикл приёмки --- от первичного разбора ответа до подтверждения готовности финальной программы --- вёлся через Codex. Поэтому результаты характеризуют работу моделей в условиях последовательного инструментального контроля и обратной связи со стороны Codex.

\section{Построение исправляющего промта}
Исправление формулируется на уровне наблюдаемого контракта. Для ошибки сохранения недостаточно потребовать «правильный NPZ»: в запросе DeepSeek указан загрузчик SciPy и требование использовать совместимый способ записи. Для ошибки словаря Алисы приведено конкретное расхождение: четыре строки вместо шести, поскольку изолированные слова не попадали в матрицу.

Промт Королёву после второй версии, напротив, разделяет две независимые проблемы: обратные слеши в путях Windows и успешное завершение при пропуске повреждённого документа. Такой запрос помогает проверить исправления отдельно. Однако эффективность этой формы запроса здесь подтверждается последовательностью результатов, а не экспериментом сравнения разных формулировок одного и того же промта.

\section{Проектирование набора проверок}
Пять наборов имеют заранее заданные тексты: базовые документы, пустые документы, пустой каталог, повторяющиеся слова и смешанные форматы. Для них независимо вычисляются ожидаемые признаки и значения. Два дополнительных сценария проверяют сообщение о повреждённом DOCX и обработку исходного корпуса.

Разделение необходимо для локализации ошибки. Сбой только на смешанных форматах указывает на чтение или преобразование текста, а несовпадение синтаксической матрицы при верной TF--IDF требует анализа зависимостей и словаря. Успешное прохождение пустого каталога показывает, что программа сохраняет согласованные структуры без вызова обучения на отсутствующем словаре.

\chapter{Технологическая часть}
\section{Выбор языка и среды программирования}
Проверка выполнялась в Python 3.12.14; для синтаксиса использовались spaCy 3.8.16 и модель \texttt{ru\_core\_news\_sm} 3.8.0. Одинаковый языковой анализатор снижает вероятность того, что различия матриц вызваны разными деревьями разбора, а не логикой программ. При этом проверяется согласованность с выбранным анализатором, а не лингвистическая точность spaCy относительно ручной разметки.

Каждая программа запускается отдельным процессом с параметрами входного и выходного каталогов. Перед запуском выполняется компиляция Python. Протокол фиксирует код завершения, длительность и хеш исходного файла, связывая измерение с определённой версией. Имена выходных файлов при необходимости сопоставляются через адаптер: он меняет адрес чтения артефакта, но не его содержимое.

\section{Реализация программного обеспечения}
Слова выделяются как последовательности русских и латинских букв и приводятся к нижнему регистру. Для каждого уникального слова считаются символьные биграммы и триграммы с границами \texttt{\^} и \texttt{\$}. Строка такой матрицы соответствует слову, поэтому повторение слова в корпусе не должно само по себе многократно увеличивать его символьное представление.

\begin{samepage}
Документная TF--IDF строится по словным униграммам и биграммам. Сглаженную обратную частоту вычисляют по формуле~\eqref{eq:idf}:
\begin{equation}\label{eq:idf}
 \mathrm{idf}(t)=\ln\frac{1+N}{1+\mathrm{df}(t)}+1,
\end{equation}

где $\mathrm{idf}(t)$~--- сглаженная обратная частота термина $t$;\\
$t$~--- термин;\\
$N$~--- число документов;\\
$\mathrm{df}(t)$~--- число документов, содержащих термин $t$.\par
\end{samepage}

\begin{samepage}
Нормированный вес термина в документе вычисляют по формуле~\eqref{eq:tfidf}:
\begin{equation}\label{eq:tfidf}
 x_{dt}=\frac{\mathrm{tf}(t,d)\,\mathrm{idf}(t)}
 {\sqrt{\sum_u(\mathrm{tf}(u,d)\mathrm{idf}(u))^2}},
\end{equation}

где $x_{dt}$~--- нормированный вес термина $t$ в документе $d$;\\
$d$~--- документ;\\
$t$~--- термин;\\
$\mathrm{tf}(t,d)$~--- частота термина $t$ в документе $d$;\\
$\mathrm{idf}(t)$~--- сглаженная обратная частота термина $t$ по формуле~\eqref{eq:idf};\\
$u$~--- индекс термина при суммировании по словарю.\par
\end{samepage}

При нулевом словаре деление не выполняется. Нулевой документ сохраняет свою строку: исключение документа ради обхода ошибки исказило бы соответствие между корпусом и матрицей.

В синтаксической матрице слов признаки OUT относятся к голове связи и содержат лемму зависимого токена; IN относятся к зависимому и содержат лемму головы. Корневые петли и пунктуация исключаются. В документной матрице накапливаются сочетания части речи головы, типа связи и части речи зависимого. Общий словарь слов формируется до отбора связей, чтобы слова без подходящего контекста оставались нулевыми строками.

\subsection{Решение GigaChat}
Версия 04 принята после устранения ошибок запуска и чтения документов. Последний корректирующий запрос указывал на отсутствие импорта \texttt{re}, неверные обращения к модулям ODT и повторный обход таблиц. Предложенный обход абзацев ODT включает текст ячеек и не требует второго накопления того же содержимого.

Дублирование текста таблицы особенно опасно тем, что не обязательно вызывает исключение. Оно меняет частоты, а затем веса и синтаксические счётчики. Поэтому для этой реализации проверка численных результатов на смешанных форматах имеет большее значение, чем сам факт успешного чтения файла.

\subsection{Решение Qwen}
Финальная версия 03 корректно разделяет извлечение текста и его использование разными алгоритмами. Для лексических признаков допустима заданная токенизация, но синтаксическому анализатору нужен исходный текст с границами предложений. Предварительная очистка, удобная для TF--IDF, не должна автоматически применяться перед dependency parsing.

Важна также отдельная обработка отсутствующего словаря. Передача пустых данных в обычный путь обучения векторизатора приводит к ошибке; корректное решение возвращает матрицы требуемой формы. Финальный PASS подтверждает согласованность этих ветвей на контрольных наборах, включая пустой каталог.

\subsection{Решение Яндекс.Алисы}
В первой доступной реализации лексические матрицы и документные синтаксические признаки базового сценария были корректны, но словная синтаксическая матрица имела неполный состав строк. Исправление состояло в предварительном построении общего множества слов, независимо от того, найдена ли для них допустимая связь.

Дополнительно потребовалось чтение TXT в CP1251 после неудачного UTF-8. Без него обрабатывались 49 документов вместо 50. Это не косметическое отклонение: меняются состав корпуса и документные частоты, поэтому программа не должна сообщать о полном успехе после исключения нечитаемого файла.

\subsection{Решение ассистента МГТУ}
Первая версия использовала несуществующую функцию ODT \texttt{get\_string}; исправляющий запрос требовал \texttt{extractText}. Другие изменения обеспечили создание всех матриц для пустых входов и поддержку CP1251. Базовая вычислительная логика при этом уже проходила проверку, что позволило сосредоточить запрос на инфраструктуре чтения.

Во второй версии осталась сериализация относительных путей через платформозависимое строковое представление. Переход к POSIX-путям с прямым слешем сделал метки переносимыми и совпадающими с манифестом корпуса. Отдельно исправлено поведение на повреждённом файле: диагностическое сообщение должно быть доступно проверке, а не теряться при успешном завершении процесса.

\subsection{Решение DeepSeek}
В версии 01 массивы разреженной матрицы сохранялись общим архиватором NumPy. Проверка использовала \texttt{scipy.sparse.load\_npz}, поэтому такой результат не проходил контракт формата. Переход к \texttt{save\_npz} устранил несовместимость без изменения смысла вычисляемых признаков.

Версия 02 ещё обращалась к классам таблиц через неверный модуль ODT. В версии 03 использованы импорты \texttt{TableRow} и \texttt{TableCell} из \texttt{odf.table}. Локальность последнего дефекта подтверждается тем, что до его исправления шесть остальных сценариев уже проходили, а после исправления прошли все семь.

\section{Протоколы взаимодействия с БЯМ}
Промты и ответы по моделям собраны в приложениях~\ref{app:gigachat} (GigaChat), \ref{app:qwen} (Qwen), \ref{app:alice} (Алиса), \ref{app:korolev} (Королёв) и~\ref{app:deepseek} (DeepSeek) в порядке версий. Для интерпретации цепочки важно связывать сообщение об ошибке с предшествующей программой, а ответ на исправляющий запрос — с последующей. Номер ответа не следует смешивать с числом повторных независимых генераций: здесь ответы продолжали один диалог.

\subsection{Протокол GigaChat}
Четыре версии отражают три перехода после первоначального ответа. Версия 02 остановилась на компиляции, поэтому её нельзя включать в сравнение времени обработки корпуса: обработки не было. Версия 03 также не прошла рабочие сценарии, и последний запрос потребовал проверить импорты и путь чтения ODT.

Этот протокол демонстрирует риск регрессии при выдаче полного файла. Даже если ответ содержит нужный алгоритм, пропущенный импорт делает его недоступным. Запрос на полную программу удобен для воспроизводимости, но требует повторной проверки всего результата, а не только обсуждавшегося фрагмента.

\subsection{Протокол Qwen}
Три версии показывают, что устранение одной ошибки библиотеки не гарантирует выполнения математического контракта. Между первыми двумя версиями итоговое число пройденных сценариев осталось прежним: проходила только диагностика повреждённого документа. Поэтому оценка исправления исключительно по исчезновению traceback была бы недостаточной.

В последнем переходе совместно проверялись пустые входы, неизменность текста для spaCy и чтение корпуса. Финальная проверка полезна именно как совокупность сценариев: она предотвращает принятие локального исправления, которое ещё не решает задачу целиком.

\subsection{Протокол Яндекс.Алисы}
Первая попытка хранится как отсутствие ответа, затем следуют две реализации. На версии 02 проверяющий запрос уже мог указать точный дефект словаря: правильные значения для части матриц сочетались с отсутствующими строками другой. Это позволило описать требуемое исправление через состав выходного результата.

После одного корректирующего перехода между программами получен PASS всех сценариев. В данной цепочке это меньше переходов между версиями кода, чем у остальных систем. Однако из этого не следует превосходство по времени или стоимости: пустая попытка тоже потребовала взаимодействия, а время ответов не измерялось.

\subsection{Протокол ассистента МГТУ}
Первый корректирующий промт объединяет проблемы ODT, кодировки, пустых входов и извлечения HTML. Второй уже уточняет переносимость путей и сообщение о повреждённом DOCX. Содержание запросов показывает переход от функциональных дефектов к согласованности внешнего интерфейса.

Итог второй версии полезен для диагностики: пять эталонных наборов прошли, но программа ещё не удовлетворяла общему критерию приёмки. Тем самым протокол подтверждает необходимость отдельных проверок эксплуатации программы поверх проверки формул.

\subsection{Протокол DeepSeek}
Первый исправляющий запрос потребовал совместимого CSR-файла и корректного поведения на пустом входе. После него общий результат изменился с нуля пройденных сценариев до шести. Второй запрос был уже узким и касался конкретного импорта ODT.

Это пример уменьшения области исправления по мере приближения к приёмке. Такой ход взаимодействия подтверждён данным диалогом, но не является статистическим доказательством, что DeepSeek всегда лучше сохраняет исправления: для такого вывода нужны повторные серии.

\chapter{Исследовательский раздел}
\section{Объекты исследования и методика оценки}
Единицей исследования является сохранённая версия программы. Для каждой доступной реализации фиксируются компиляция, результат сценариев и время отдельного процесса. Итоговые машинные отчёты приведены в приложении~\ref{app:json}, сводка — в приложении~\ref{app:summary}. Результаты GigaChat приведены в разделе~\ref{app:json:gigachat}, Qwen — в разделе~\ref{app:json:qwen}, Алисы — в разделе~\ref{app:json:alice}, Королёва — в разделе~\ref{app:json:korolev}, DeepSeek — в разделе~\ref{app:json:deepseek}. Таким образом, читатель может проверить выводы по документу без доступа к рабочему каталогу.

Для пяти контрольных наборов эталон независимо вычисляет четыре матрицы. Проверяются упорядоченные уникальные метки, размеры, CSR-тип и конечность неотрицательных значений. Сравнение чисел выполняется при совпадении меток и формы с допустимым максимальным отклонением менее $10^{-6}$. Для TF--IDF дополнительно проверяется единичная либо нулевая норма строки, для счётчиков — целочисленность.

Полный корпус проверяется иначе: контролируются структурные свойства и точное наличие всех 50 путей в документных матрицах. Независимый числовой эталон для всего корпуса не вычисляется. Наличие metadata проверяется как наличие файла, а не как полная валидация каждого его поля. На повреждённом DOCX принимается журнал с именем файла либо ненулевое завершение с диагностикой; продолжение обработки остальных файлов этим сценарием не доказывается.

\section{Ход исследования и анализ взаимодействия}
\subsection{GigaChat}
Версия 01 прошла только диагностику повреждённого файла: 1 из 7 сценариев. Версия 02 не скомпилировалась, версия 03 получила 0 из 7, версия 04 — 7 из 7. Последовательность не является монотонным улучшением: промежуточные ответы ухудшали доступность результата.

Следовательно, преимущество итоговой программы следует отделять от затрат на её получение. Финальный результат принят, но в этой цепочке контроль импортов и компиляции мог раньше выявить проблемы и сделать обратную связь короче. Большое число строк ответа само по себе не служит положительным критерием качества.

\subsection{Qwen}
Версии 01 и 02 прошли по 1 из 7 сценариев, версия 03 — все семь. Ранний успех в диагностике повреждённого файла не сопровождался корректным вычислительным результатом. Два одинаковых агрегированных показателя также не означают идентичные программы: состав внутренних дефектов мог меняться.

Основной вывод состоит в необходимости проверять несколько уровней результата. Исправление вызова библиотеки снимает ошибку выполнения, но затем становятся видны несоответствия словаря или синтаксических признаков. Оценка только по количеству исключений скрыла бы этот переход.

\subsection{Яндекс.Алиса}
Попытка 01 не содержит кода; версия 02 прошла 2 из 7 сценариев — пустые документы и пустой каталог. Версия 03 прошла 7 из 7. Следовательно, утверждать, что обе доступные программы полностью работали, нельзя: первая только частично удовлетворяла проверкам.

Исправление состава словаря и кодировки устранило ошибки, которые не были очевидны по виду выходных файлов. Наличие NPZ и успешное завершение процесса не обеспечивали полноты результата. Для этой модели решающим оказался контроль меток и числа документов.

\subsection{Ассистент МГТУ}
Версия 01 прошла 3 из 6 выполненных сценариев; полного корпусного запуска в её JSON нет. Версия 02 прошла 5 из 7, версия 03 — 7 из 7. Разные знаменатели нужно сохранять: запись 3 из 7 для первой версии ошибочно представляла бы невыполненный тест как проваленный.

По базовой и повторной лексике первая версия Королёва была ближе к требуемому результату, чем первые реализации других участников. При этом итоговая приёмка всё равно потребовала двух исправлений. Этот случай показывает, что хорошая математическая основа и готовность программы как воспроизводимого инструмента — разные критерии.

\subsection{DeepSeek}
Результаты последовательных версий составили 0 из 7, 6 из 7 и 7 из 7. Нулевой старт здесь во многом объясняется общим барьером сериализации: загрузчик не мог прочитать матрицы и перейти к проверке их чисел. Поэтому он не означает, что каждая математическая формула была ошибочной.

Особенно показателен проход полного корпуса во второй версии при провале formats. Специально сконструированный тест обнаружил ошибку, которую большой набор не выявил. Это обосновывает сочетание реальных документов и малых направленных примеров в методике оценки.

\section{Сравнительный анализ алгоритмов}
Сопоставление разделяет корректность расчёта, совместимость сохранения и устойчивость чтения. По начальной корректности двух содержательных эталонных сценариев выделяется Королёв. По числу переходов между доступными программами до полного PASS выделяется Алиса: один переход, но с предшествующей пустой попыткой. По итоговой проверяемой функциональности победителя нет — все пять финальных решений приняты.

Алгоритмы используют общую заданную схему признаков и один синтаксический анализатор. Поэтому итоговое совпадение подтверждает реализацию спецификации, а не открытие разных по качеству методов обработки языка. Различия преимущественно находятся в извлечении текста, построении словарей, обработке исключений и сохранении результатов.

Сведения о прохождении сценариев, числе версий и времени обработки корпуса приведены в таблице~\ref{tab:chains}.

\begin{table}[H]
\centering\small
\caption{Результаты сохранённых цепочек}
\label{tab:chains}
\begin{tabular}{|p{2.45cm}|p{4.5cm}|>{\centering\arraybackslash}p{2.85cm}|>{\centering\arraybackslash}p{3.25cm}|}
\hline
\centering Модель & \centering Проход сценариев по версиям & Программных версий & Полный корпус, с \\\hline
GigaChat & 1/7; ошибка компиляции; 0/7; 7/7 & 4 & 5,678 \\\hline
Qwen & 1/7; 1/7; 7/7 & 3 & 6,373 \\\hline
Алиса & нет кода; 2/7; 7/7 & 2 & 6,070 \\\hline
Королёв & 3/6; 5/7; 7/7 & 3 & 6,259 \\\hline
DeepSeek & 0/7; 6/7; 7/7 & 3 & 10,021 \\\hline
\end{tabular}
\end{table}

В последнем столбце приведено время процесса финальной программы на корпусе из её JSON. Это однократные измерения с загрузкой зависимостей и модели, а не средняя скорость чистого алгоритма. GigaChat имеет наименьшее зафиксированное время, DeepSeek — наибольшее; разница примерно 1,76 раза. Повторных прогонов для оценки разброса нет, поэтому использовать эти числа как устойчивый рейтинг производительности нельзя.

Стоимость промтинга характеризуется числом программ и попыток, но не токенами или секундами генерации: последние не измерены. Проход 7/7 означает успешность семи различных сценариев одной программы, а не вероятность успеха модели 100 процентов при новом запросе. Это различие предотвращает завышенную интерпретацию результата.

\section{Выводы по исследовательскому разделу}
Для данного задания работоспособный результат получен от всех участников, но ни одна первая попытка не была достаточна для полной приёмки. Королёв лучше стартовал на базовой вычислительной логике; у Алисы было меньше переходов между программами после получения кода; DeepSeek показал наиболее явный переход от общего дефекта формата к одному локальному дефекту чтения. GigaChat потребовал больше версий, несмотря на минимальное наблюдаемое время финального запуска.

Эти преимущества относятся к разным осям сравнения и не образуют единого победителя. Практический результат исследования — необходимость заранее определить критерий выбора: готовность первой реализации, число корректирующих ответов, полнота обработки либо время исполнения. Объединение этих показателей в субъективный балл без весов исказило бы вывод.

\section{Стабильность исправлений и ограничения исследования}
\subsection{Методика анализа стабильности}
В данном эксперименте стабильность понимается как сохранение пройденных проверок при переходе к новой версии. Это отличается от повторяемости независимой генерации одного промта. В частности, ухудшение результатов GigaChat на промежуточном шаге показывает, что новый полный ответ необходимо проверять целиком.

Для строгого исследования вероятности успеха потребовались бы несколько независимых диалогов на каждую модель, одинаковые ограничения контекста и учёт точных версий сервиса. Здесь сохранена одна цепочка на участника, поэтому можно анализировать наблюдаемые изменения, но нельзя оценивать частоту аналогичных ошибок в генеральной совокупности ответов.

\subsection{Влияние режима рассуждений}
Контролируемой серии с включённым и выключенным режимом рассуждений не проводилось. Наличие такого раздела в образце отчёта не создаёт соответствующих экспериментальных данных для этой работы. Поэтому связывать найденные ошибки с конкретной архитектурой рассуждения или утверждать преимущество одного режима было бы необоснованно.

Возможное продолжение исследования — повторить один и тот же начальный запрос в фиксированных режимах, не подмешивая разные исправляющие подсказки. Тогда отдельно измеряются доля корректных первых программ, число необходимых уточнений и вариативность результатов; текущая работа предоставляет методику проверки для такой серии.

\subsection{Сводные результаты и границы применимости}
Пять принятых финальных решений подтверждают достижимость задачи через итеративный промтинг. Согласованность контрольных матриц и охват корпуса дополняют друг друга: первое проверяет заданную логику на известных данных, второе — способность обработать реальный набор и сохранить согласованный результат.

Остаются непроверенными все возможные разновидности вложенных таблиц, смешанные кодировки внутри файла и лингвистическая точность зависимостей. Полученный результат достаточен для описанного набора критериев; перенос на другой корпус требует новых направленных тестов, особенно для форматов, где ранее наблюдались ошибки.

\chapter*{ЗАКЛЮЧЕНИЕ}
\addcontentsline{toc}{chapter}{ЗАКЛЮЧЕНИЕ}
В работе получены и проверены программы пяти БЯМ для четырёх способов представления текста. Сравнение показало, что формальное сходство ответов с полноценной программой не гарантирует результата: ранние версии нарушали контракт сохранения, теряли документы или строки словаря, а в одном случае не компилировались. Проверка только запуска не выявила бы все эти дефекты.

По конечной функциональности участники равны: принятые версии прошли семь сценариев, включая корпус из 50 документов. Однако одинаковый конечный результат скрывает различия в объёме необходимого контроля и характере помощи со стороны пользователя. Поэтому аналитический вывод по моделям следует связывать с конкретным этапом получения программы.

\emph{Ассистент МГТУ (Королёв)} показал наиболее убедительную исходную реализацию базовой вычислительной логики: первая программа уже проходила базовый сценарий и повторяющиеся слова. Последующие запросы в основном доводили обработку форматов, пустых входов и внешний контракт результата. В рамках этой задачи его преимущество состоит в меньшей необходимости уточнять основные вычисления, хотя готовность к эксплуатации всё равно потребовала двух исправлений.

\emph{Яндекс.Алиса} потребовала только одного корректирующего перехода между двумя доступными программами. Это наиболее короткий цикл исправления кода среди исследованных цепочек, но ему предшествовала попытка без ответа. Следовательно, её результат выгоден по числу исправлений после получения программы, однако не позволяет объявить её лучшей по надёжности предоставления первого ответа. Ошибки словаря и кодировки также показывают, что внешне работающая программа нуждалась в проверке полноты данных.

\emph{Qwen} достиг полного PASS в третьей версии, причём первые две проходили только диагностику повреждённого файла. Его результат подчёркивает различие между наличием защитных ветвей программы и правильностью основного алгоритма. Для этой цепочки решающим было уточнение обработки пустых входов и подготовки текста к синтаксическому разбору; по ранней готовности решения Qwen уступил Королёву, а по итоговой проверяемой функциональности сравнялся с ним.

\emph{DeepSeek} продемонстрировал наиболее отчётливое сужение области оставшихся дефектов: после устранения общей несовместимости NPZ вторая версия прошла шесть сценариев из семи, а последний запрос исправил конкретный импорт ODT. Такая динамика удобна для адресного промтинга: становится понятно, какой компонент ещё мешает приёмке. При этом начальный дефект сериализации был критичен для использования результата, а финальная программа имела наибольшее время единственного измеренного корпусного запуска.

\emph{GigaChat} потребовал четыре версии и показал промежуточные ошибки компиляции и запуска. В данной цепочке это означает наибольшее число программных ответов до приёмки и необходимость особенно внимательно проверять целостность нового полного файла. Одновременно его финальная программа имела наименьшее зафиксированное время обработки корпуса. Таким образом, затраты на получение программы и скорость её исполнения не совпадают: более длинный диалог не обязательно приводит к более медленному коду.

Совокупный вывод состоит в том, что для этой задачи Королёв выделяется исходной готовностью вычислений, Алиса — коротким исправлением после получения кода, DeepSeek — хорошо локализуемым последним дефектом, а GigaChat — временем финального запуска при более длинной отладке. Qwen достиг того же функционального результата после уточнения нескольких взаимосвязанных требований. Эти наблюдения дают критерии выбора и контроля, но не универсальный рейтинг: один диалог на модель не устанавливает устойчивого превосходства. Общим условием успешного промтинга для всех участников оказался независимый контроль результата после каждого ответа.

Наиболее содержательный результат анализа промтинга заключается в роли проверяемой обратной связи. Запрос с конкретным расхождением переводит общую задачу в локальное требование: сохранить нулевые строки, прочитать CP1251, записать SciPy-совместимый NPZ или исправить модуль импорта. В наблюдаемых цепочках такие уточнения завершились успешной приёмкой. Сравнения с альтернативными стратегиями промтинга не проводилось, поэтому это результат применённого процесса, а не доказательство оптимальности формулировок.

Практически получение программы следует организовывать как цикл с независимым критерием завершения. Сначала задаются формат результата и контрольные примеры; затем запускается неизменённый ответ, выявляется причина расхождения и формулируется исправляющий запрос. Новая версия проходит весь применимый набор повторно. Этот порядок позволяет использовать БЯМ как генератор реализации, сохраняя проверку правильности вне модели.

Выводы ограничены одной задачей и одной цепочкой на сервис. Точные числовые эталоны использованы на контрольных наборах; для большого корпуса проверены структурные свойства и охват документов. Следующий обоснованный шаг — независимые повторные генерации и многократные измерения времени. Они позволят оценить устойчивость уже статистически, тогда как настоящий отчёт документирует конкретный воспроизводимый путь к рабочим программам.


\begin{thebibliography}{9}
\addcontentsline{toc}{chapter}{СПИСОК ИСПОЛЬЗОВАННЫХ ИСТОЧНИКОВ}
\bibitem{scipy} Sparse arrays (scipy.sparse) // SciPy Manual.~--- URL: \url{https://docs.scipy.org/doc/scipy/reference/sparse.html} (дата обращения: 17.11.2026).
\bibitem{sklearn} TfidfTransformer // scikit-learn Documentation.~--- URL: \url{https://scikit-learn.org/stable/modules/generated/sklearn.feature_extraction.text.TfidfTransformer.html} (дата обращения: 17.11.2026).
\bibitem{spacy} Linguistic Features // spaCy Usage Documentation.~--- URL: \url{https://spacy.io/usage/linguistic-features} (дата обращения: 17.11.2026).
\bibitem{gigachat} Руководства GigaChat API // Сбер : документация для разработчиков.~--- URL: \url{https://developers.sber.ru/docs/ru/gigachat/guides/main} (дата обращения: 17.11.2026).
\bibitem{qwen} Qwen : официальный репозиторий // GitHub.~--- URL: \url{https://github.com/QwenLM/Qwen} (дата обращения: 17.11.2026).
\bibitem{alice} Алиса AI // Яндекс : сайт.~--- URL: \url{https://alice.yandex.ru/} (дата обращения: 17.11.2026).
\bibitem{korolev} Королёв ИИ : научно-образовательная платформа // МГТУ им. Н.Э. Баумана.~--- URL: \url{https://korolev-ai.bmstu.ru/} (дата обращения: 17.11.2026).
\bibitem{deepseek} DeepSeek : официальный сайт.~--- URL: \url{https://www.deepseek.com/} (дата обращения: 17.11.2026).
\bibitem{codex} Codex // OpenAI : сайт.~--- URL: \url{https://openai.com/codex/} (дата обращения: 17.11.2026).
\end{thebibliography}




\BeginAppendices
\prilojenie[обязательное]{Промты и ответы: GigaChat}
\label{app:gigachat}
Ниже приведены исходные тексты без сокращений и редакторских замен. Служебные маркеры и Markdown сохранены как часть оригинала. Листинги продолжаются на следующих страницах без пропуска строк.
\inputcode[style=logstyle,firstline=1,lastline=12,firstnumber=1,caption={GigaChat: промт версии 1}]{../prompts/gigachat_v1_prompt.txt}
\inputcode[style=logstyle,firstline=1,lastline=80,firstnumber=1,caption={GigaChat: ответ версии 1}]{../prompts/gigachat_v1_response.txt}
\inputcode[style=logstyle,firstline=81,lastline=160,firstnumber=81,caption={GigaChat: ответ версии 1}]{../prompts/gigachat_v1_response.txt}
\inputcode[style=logstyle,firstline=161,lastline=240,firstnumber=161,caption={GigaChat: ответ версии 1}]{../prompts/gigachat_v1_response.txt}
\inputcode[style=logstyle,firstline=241,lastline=320,firstnumber=241,caption={GigaChat: ответ версии 1}]{../prompts/gigachat_v1_response.txt}
\inputcode[style=logstyle,firstline=321,lastline=400,firstnumber=321,caption={GigaChat: ответ версии 1}]{../prompts/gigachat_v1_response.txt}
\inputcode[style=logstyle,firstline=401,lastline=480,firstnumber=401,caption={GigaChat: ответ версии 1}]{../prompts/gigachat_v1_response.txt}
\inputcode[style=logstyle,firstline=481,lastline=560,firstnumber=481,caption={GigaChat: ответ версии 1}]{../prompts/gigachat_v1_response.txt}
\inputcode[style=logstyle,firstline=561,lastline=594,firstnumber=561,caption={GigaChat: ответ версии 1}]{../prompts/gigachat_v1_response.txt}
\inputcode[style=logstyle,firstline=1,lastline=1,firstnumber=1,caption={GigaChat: промт версии 2}]{../prompts/gigachat_v2_prompt.txt}
\inputcode[style=logstyle,firstline=1,lastline=80,firstnumber=1,caption={GigaChat: ответ версии 2}]{../prompts/gigachat_v2_response.txt}
\inputcode[style=logstyle,firstline=81,lastline=160,firstnumber=81,caption={GigaChat: ответ версии 2}]{../prompts/gigachat_v2_response.txt}
\inputcode[style=logstyle,firstline=161,lastline=240,firstnumber=161,caption={GigaChat: ответ версии 2}]{../prompts/gigachat_v2_response.txt}
\inputcode[style=logstyle,firstline=241,lastline=320,firstnumber=241,caption={GigaChat: ответ версии 2}]{../prompts/gigachat_v2_response.txt}
\inputcode[style=logstyle,firstline=321,lastline=400,firstnumber=321,caption={GigaChat: ответ версии 2}]{../prompts/gigachat_v2_response.txt}
\inputcode[style=logstyle,firstline=401,lastline=443,firstnumber=401,caption={GigaChat: ответ версии 2}]{../prompts/gigachat_v2_response.txt}
\inputcode[style=logstyle,firstline=1,lastline=1,firstnumber=1,caption={GigaChat: промт версии 3}]{../prompts/gigachat_v3_prompt.txt}
\inputcode[style=logstyle,firstline=1,lastline=80,firstnumber=1,caption={GigaChat: ответ версии 3}]{../prompts/gigachat_v3_response.txt}
\inputcode[style=logstyle,firstline=81,lastline=160,firstnumber=81,caption={GigaChat: ответ версии 3}]{../prompts/gigachat_v3_response.txt}
\inputcode[style=logstyle,firstline=161,lastline=240,firstnumber=161,caption={GigaChat: ответ версии 3}]{../prompts/gigachat_v3_response.txt}
\inputcode[style=logstyle,firstline=241,lastline=320,firstnumber=241,caption={GigaChat: ответ версии 3}]{../prompts/gigachat_v3_response.txt}
\inputcode[style=logstyle,firstline=321,lastline=335,firstnumber=321,caption={GigaChat: ответ версии 3}]{../prompts/gigachat_v3_response.txt}
\inputcode[style=logstyle,firstline=1,lastline=1,firstnumber=1,caption={GigaChat: промт версии 4}]{../prompts/gigachat_v4_prompt.txt}
\inputcode[style=logstyle,firstline=1,lastline=80,firstnumber=1,caption={GigaChat: ответ версии 4}]{../prompts/gigachat_v4_response.txt}
\inputcode[style=logstyle,firstline=81,lastline=160,firstnumber=81,caption={GigaChat: ответ версии 4}]{../prompts/gigachat_v4_response.txt}
\inputcode[style=logstyle,firstline=161,lastline=240,firstnumber=161,caption={GigaChat: ответ версии 4}]{../prompts/gigachat_v4_response.txt}
\inputcode[style=logstyle,firstline=241,lastline=320,firstnumber=241,caption={GigaChat: ответ версии 4}]{../prompts/gigachat_v4_response.txt}
\inputcode[style=logstyle,firstline=321,lastline=324,firstnumber=321,caption={GigaChat: ответ версии 4}]{../prompts/gigachat_v4_response.txt}
\prilojenie[обязательное]{Промты и ответы: Qwen}
\label{app:qwen}
Ниже приведены исходные тексты без сокращений и редакторских замен. Служебные маркеры и Markdown сохранены как часть оригинала. Листинги продолжаются на следующих страницах без пропуска строк.
\inputcode[style=logstyle,firstline=1,lastline=12,firstnumber=1,caption={Qwen: промт версии 1}]{../prompts/qwen_v1_prompt.txt}
\inputcode[style=logstyle,firstline=1,lastline=80,firstnumber=1,caption={Qwen: ответ версии 1}]{../prompts/qwen_v1_response.txt}
\inputcode[style=logstyle,firstline=81,lastline=160,firstnumber=81,caption={Qwen: ответ версии 1}]{../prompts/qwen_v1_response.txt}
\inputcode[style=logstyle,firstline=161,lastline=240,firstnumber=161,caption={Qwen: ответ версии 1}]{../prompts/qwen_v1_response.txt}
\inputcode[style=logstyle,firstline=241,lastline=320,firstnumber=241,caption={Qwen: ответ версии 1}]{../prompts/qwen_v1_response.txt}
\inputcode[style=logstyle,firstline=321,lastline=400,firstnumber=321,caption={Qwen: ответ версии 1}]{../prompts/qwen_v1_response.txt}
\inputcode[style=logstyle,firstline=401,lastline=480,firstnumber=401,caption={Qwen: ответ версии 1}]{../prompts/qwen_v1_response.txt}
\inputcode[style=logstyle,firstline=481,lastline=560,firstnumber=481,caption={Qwen: ответ версии 1}]{../prompts/qwen_v1_response.txt}
\inputcode[style=logstyle,firstline=561,lastline=606,firstnumber=561,caption={Qwen: ответ версии 1}]{../prompts/qwen_v1_response.txt}
\inputcode[style=logstyle,firstline=1,lastline=1,firstnumber=1,caption={Qwen: промт версии 2}]{../prompts/qwen_v2_prompt.txt}
\inputcode[style=logstyle,firstline=1,lastline=80,firstnumber=1,caption={Qwen: ответ версии 2}]{../prompts/qwen_v2_response.txt}
\inputcode[style=logstyle,firstline=81,lastline=160,firstnumber=81,caption={Qwen: ответ версии 2}]{../prompts/qwen_v2_response.txt}
\inputcode[style=logstyle,firstline=161,lastline=240,firstnumber=161,caption={Qwen: ответ версии 2}]{../prompts/qwen_v2_response.txt}
\inputcode[style=logstyle,firstline=241,lastline=320,firstnumber=241,caption={Qwen: ответ версии 2}]{../prompts/qwen_v2_response.txt}
\inputcode[style=logstyle,firstline=321,lastline=400,firstnumber=321,caption={Qwen: ответ версии 2}]{../prompts/qwen_v2_response.txt}
\inputcode[style=logstyle,firstline=401,lastline=413,firstnumber=401,caption={Qwen: ответ версии 2}]{../prompts/qwen_v2_response.txt}
\inputcode[style=logstyle,firstline=1,lastline=1,firstnumber=1,caption={Qwen: промт версии 3}]{../prompts/qwen_v3_prompt.txt}
\inputcode[style=logstyle,firstline=1,lastline=80,firstnumber=1,caption={Qwen: ответ версии 3}]{../prompts/qwen_v3_response.txt}
\inputcode[style=logstyle,firstline=81,lastline=160,firstnumber=81,caption={Qwen: ответ версии 3}]{../prompts/qwen_v3_response.txt}
\inputcode[style=logstyle,firstline=161,lastline=181,firstnumber=161,caption={Qwen: ответ версии 3}]{../prompts/qwen_v3_response.txt}
\prilojenie[обязательное]{Промты и ответы: Яндекс.Алиса}
\label{app:alice}
Ниже приведены исходные тексты без сокращений и редакторских замен. Служебные маркеры и Markdown сохранены как часть оригинала. Листинги продолжаются на следующих страницах без пропуска строк.
\inputcode[style=logstyle,firstline=1,lastline=12,firstnumber=1,caption={Яндекс.Алиса: промт версии 1}]{../prompts/alice_v1_prompt.txt}
\inputcode[style=logstyle,firstline=1,lastline=1,firstnumber=1,caption={Яндекс.Алиса: статус попытки 1}]{../prompts/alice_v1_status.txt}
\inputcode[style=logstyle,firstline=1,lastline=1,firstnumber=1,caption={Яндекс.Алиса: промт версии 2}]{../prompts/alice_v2_prompt.txt}
\inputcode[style=logstyle,firstline=1,lastline=80,firstnumber=1,caption={Яндекс.Алиса: ответ версии 2}]{../prompts/alice_v2_response.txt}
\inputcode[style=logstyle,firstline=81,lastline=160,firstnumber=81,caption={Яндекс.Алиса: ответ версии 2}]{../prompts/alice_v2_response.txt}
\inputcode[style=logstyle,firstline=161,lastline=240,firstnumber=161,caption={Яндекс.Алиса: ответ версии 2}]{../prompts/alice_v2_response.txt}
\inputcode[style=logstyle,firstline=241,lastline=320,firstnumber=241,caption={Яндекс.Алиса: ответ версии 2}]{../prompts/alice_v2_response.txt}
\inputcode[style=logstyle,firstline=321,lastline=400,firstnumber=321,caption={Яндекс.Алиса: ответ версии 2}]{../prompts/alice_v2_response.txt}
\inputcode[style=logstyle,firstline=401,lastline=480,firstnumber=401,caption={Яндекс.Алиса: ответ версии 2}]{../prompts/alice_v2_response.txt}
\inputcode[style=logstyle,firstline=481,lastline=560,firstnumber=481,caption={Яндекс.Алиса: ответ версии 2}]{../prompts/alice_v2_response.txt}
\inputcode[style=logstyle,firstline=561,lastline=565,firstnumber=561,caption={Яндекс.Алиса: ответ версии 2}]{../prompts/alice_v2_response.txt}
\inputcode[style=logstyle,firstline=1,lastline=1,firstnumber=1,caption={Яндекс.Алиса: промт версии 3}]{../prompts/alice_v3_prompt.txt}
\inputcode[style=logstyle,firstline=1,lastline=80,firstnumber=1,caption={Яндекс.Алиса: ответ версии 3}]{../prompts/alice_v3_response.txt}
\inputcode[style=logstyle,firstline=81,lastline=160,firstnumber=81,caption={Яндекс.Алиса: ответ версии 3}]{../prompts/alice_v3_response.txt}
\inputcode[style=logstyle,firstline=161,lastline=240,firstnumber=161,caption={Яндекс.Алиса: ответ версии 3}]{../prompts/alice_v3_response.txt}
\inputcode[style=logstyle,firstline=241,lastline=320,firstnumber=241,caption={Яндекс.Алиса: ответ версии 3}]{../prompts/alice_v3_response.txt}
\inputcode[style=logstyle,firstline=321,lastline=400,firstnumber=321,caption={Яндекс.Алиса: ответ версии 3}]{../prompts/alice_v3_response.txt}
\inputcode[style=logstyle,firstline=401,lastline=480,firstnumber=401,caption={Яндекс.Алиса: ответ версии 3}]{../prompts/alice_v3_response.txt}
\inputcode[style=logstyle,firstline=481,lastline=525,firstnumber=481,caption={Яндекс.Алиса: ответ версии 3}]{../prompts/alice_v3_response.txt}
\prilojenie[обязательное]{Промты и ответы: Королёв}
\label{app:korolev}
Ниже приведены исходные тексты без сокращений и редакторских замен. Служебные маркеры и Markdown сохранены как часть оригинала. Листинги продолжаются на следующих страницах без пропуска строк.
\inputcode[style=logstyle,firstline=1,lastline=13,firstnumber=1,caption={Королёв: промт версии 1}]{../prompts/korolev_v1_prompt.txt}
\inputcode[style=logstyle,firstline=1,lastline=80,firstnumber=1,caption={Королёв: ответ версии 1}]{../prompts/korolev_v1_response.txt}
\inputcode[style=logstyle,firstline=81,lastline=160,firstnumber=81,caption={Королёв: ответ версии 1}]{../prompts/korolev_v1_response.txt}
\inputcode[style=logstyle,firstline=161,lastline=240,firstnumber=161,caption={Королёв: ответ версии 1}]{../prompts/korolev_v1_response.txt}
\inputcode[style=logstyle,firstline=241,lastline=320,firstnumber=241,caption={Королёв: ответ версии 1}]{../prompts/korolev_v1_response.txt}
\inputcode[style=logstyle,firstline=321,lastline=329,firstnumber=321,caption={Королёв: ответ версии 1}]{../prompts/korolev_v1_response.txt}
\inputcode[style=logstyle,firstline=1,lastline=12,firstnumber=1,caption={Королёв: промт версии 2}]{../prompts/korolev_v2_prompt.txt}
\inputcode[style=logstyle,firstline=1,lastline=80,firstnumber=1,caption={Королёв: ответ версии 2}]{../prompts/korolev_v2_response.txt}
\inputcode[style=logstyle,firstline=81,lastline=160,firstnumber=81,caption={Королёв: ответ версии 2}]{../prompts/korolev_v2_response.txt}
\inputcode[style=logstyle,firstline=161,lastline=240,firstnumber=161,caption={Королёв: ответ версии 2}]{../prompts/korolev_v2_response.txt}
\inputcode[style=logstyle,firstline=241,lastline=245,firstnumber=241,caption={Королёв: ответ версии 2}]{../prompts/korolev_v2_response.txt}
\inputcode[style=logstyle,firstline=1,lastline=7,firstnumber=1,caption={Королёв: промт версии 3}]{../prompts/korolev_v3_prompt.txt}
\inputcode[style=logstyle,firstline=1,lastline=80,firstnumber=1,caption={Королёв: ответ версии 3}]{../prompts/korolev_v3_response.txt}
\inputcode[style=logstyle,firstline=81,lastline=160,firstnumber=81,caption={Королёв: ответ версии 3}]{../prompts/korolev_v3_response.txt}
\inputcode[style=logstyle,firstline=161,lastline=240,firstnumber=161,caption={Королёв: ответ версии 3}]{../prompts/korolev_v3_response.txt}
\inputcode[style=logstyle,firstline=241,lastline=243,firstnumber=241,caption={Королёв: ответ версии 3}]{../prompts/korolev_v3_response.txt}
\prilojenie[обязательное]{Промты и ответы: DeepSeek}
\label{app:deepseek}
Ниже приведены исходные тексты без сокращений и редакторских замен. Служебные маркеры и Markdown сохранены как часть оригинала. Листинги продолжаются на следующих страницах без пропуска строк.
\inputcode[style=logstyle,firstline=1,lastline=16,firstnumber=1,caption={DeepSeek: промт версии 1}]{../prompts/deepseek_v1_prompt.txt}
\inputcode[style=logstyle,firstline=1,lastline=80,firstnumber=1,caption={DeepSeek: ответ версии 1}]{../prompts/deepseek_v1_response.txt}
\inputcode[style=logstyle,firstline=81,lastline=160,firstnumber=81,caption={DeepSeek: ответ версии 1}]{../prompts/deepseek_v1_response.txt}
\inputcode[style=logstyle,firstline=161,lastline=240,firstnumber=161,caption={DeepSeek: ответ версии 1}]{../prompts/deepseek_v1_response.txt}
\inputcode[style=logstyle,firstline=241,lastline=320,firstnumber=241,caption={DeepSeek: ответ версии 1}]{../prompts/deepseek_v1_response.txt}
\inputcode[style=logstyle,firstline=321,lastline=400,firstnumber=321,caption={DeepSeek: ответ версии 1}]{../prompts/deepseek_v1_response.txt}
\inputcode[style=logstyle,firstline=401,lastline=480,firstnumber=401,caption={DeepSeek: ответ версии 1}]{../prompts/deepseek_v1_response.txt}
\inputcode[style=logstyle,firstline=481,lastline=560,firstnumber=481,caption={DeepSeek: ответ версии 1}]{../prompts/deepseek_v1_response.txt}
\inputcode[style=logstyle,firstline=561,lastline=640,firstnumber=561,caption={DeepSeek: ответ версии 1}]{../prompts/deepseek_v1_response.txt}
\inputcode[style=logstyle,firstline=641,lastline=678,firstnumber=641,caption={DeepSeek: ответ версии 1}]{../prompts/deepseek_v1_response.txt}
\inputcode[style=logstyle,firstline=1,lastline=17,firstnumber=1,caption={DeepSeek: промт версии 2}]{../prompts/deepseek_v2_prompt.txt}
\inputcode[style=logstyle,firstline=1,lastline=80,firstnumber=1,caption={DeepSeek: ответ версии 2}]{../prompts/deepseek_v2_response.txt}
\inputcode[style=logstyle,firstline=81,lastline=160,firstnumber=81,caption={DeepSeek: ответ версии 2}]{../prompts/deepseek_v2_response.txt}
\inputcode[style=logstyle,firstline=161,lastline=240,firstnumber=161,caption={DeepSeek: ответ версии 2}]{../prompts/deepseek_v2_response.txt}
\inputcode[style=logstyle,firstline=241,lastline=320,firstnumber=241,caption={DeepSeek: ответ версии 2}]{../prompts/deepseek_v2_response.txt}
\inputcode[style=logstyle,firstline=321,lastline=400,firstnumber=321,caption={DeepSeek: ответ версии 2}]{../prompts/deepseek_v2_response.txt}
\inputcode[style=logstyle,firstline=401,lastline=480,firstnumber=401,caption={DeepSeek: ответ версии 2}]{../prompts/deepseek_v2_response.txt}
\inputcode[style=logstyle,firstline=481,lastline=560,firstnumber=481,caption={DeepSeek: ответ версии 2}]{../prompts/deepseek_v2_response.txt}
\inputcode[style=logstyle,firstline=561,lastline=607,firstnumber=561,caption={DeepSeek: ответ версии 2}]{../prompts/deepseek_v2_response.txt}
\inputcode[style=logstyle,firstline=1,lastline=13,firstnumber=1,caption={DeepSeek: промт версии 3}]{../prompts/deepseek_v3_prompt.txt}
\inputcode[style=logstyle,firstline=1,lastline=80,firstnumber=1,caption={DeepSeek: ответ версии 3}]{../prompts/deepseek_v3_response.txt}
\inputcode[style=logstyle,firstline=81,lastline=160,firstnumber=81,caption={DeepSeek: ответ версии 3}]{../prompts/deepseek_v3_response.txt}
\inputcode[style=logstyle,firstline=161,lastline=240,firstnumber=161,caption={DeepSeek: ответ версии 3}]{../prompts/deepseek_v3_response.txt}
\inputcode[style=logstyle,firstline=241,lastline=320,firstnumber=241,caption={DeepSeek: ответ версии 3}]{../prompts/deepseek_v3_response.txt}
\inputcode[style=logstyle,firstline=321,lastline=400,firstnumber=321,caption={DeepSeek: ответ версии 3}]{../prompts/deepseek_v3_response.txt}
\inputcode[style=logstyle,firstline=401,lastline=480,firstnumber=401,caption={DeepSeek: ответ версии 3}]{../prompts/deepseek_v3_response.txt}
\inputcode[style=logstyle,firstline=481,lastline=553,firstnumber=481,caption={DeepSeek: ответ версии 3}]{../prompts/deepseek_v3_response.txt}
\prilojenie{Итоговые JSON-отчёты}
\label{app:json}
Локальные абсолютные пути заменены на маркер LOCAL\_PATH для обезличивания. Числовые результаты, статусы и хеши программ сохранены.
\section{GigaChat}
\label{app:json:gigachat}
\inputcode[style=logstyle,firstline=1,lastline=80,firstnumber=1,caption={GigaChat: итоговый JSON}]{../prompts/eval_gigachat.json}
\inputcode[style=logstyle,firstline=81,lastline=160,firstnumber=81,caption={GigaChat: итоговый JSON}]{../prompts/eval_gigachat.json}
\inputcode[style=logstyle,firstline=161,lastline=240,firstnumber=161,caption={GigaChat: итоговый JSON}]{../prompts/eval_gigachat.json}
\inputcode[style=logstyle,firstline=241,lastline=320,firstnumber=241,caption={GigaChat: итоговый JSON}]{../prompts/eval_gigachat.json}
\inputcode[style=logstyle,firstline=321,lastline=400,firstnumber=321,caption={GigaChat: итоговый JSON}]{../prompts/eval_gigachat.json}
\inputcode[style=logstyle,firstline=401,lastline=480,firstnumber=401,caption={GigaChat: итоговый JSON}]{../prompts/eval_gigachat.json}
\inputcode[style=logstyle,firstline=481,lastline=560,firstnumber=481,caption={GigaChat: итоговый JSON}]{../prompts/eval_gigachat.json}
\inputcode[style=logstyle,firstline=561,lastline=617,firstnumber=561,caption={GigaChat: итоговый JSON}]{../prompts/eval_gigachat.json}
\section{Qwen}
\label{app:json:qwen}
\inputcode[style=logstyle,firstline=1,lastline=80,firstnumber=1,caption={Qwen: итоговый JSON}]{../prompts/eval_qwen.json}
\inputcode[style=logstyle,firstline=81,lastline=160,firstnumber=81,caption={Qwen: итоговый JSON}]{../prompts/eval_qwen.json}
\inputcode[style=logstyle,firstline=161,lastline=240,firstnumber=161,caption={Qwen: итоговый JSON}]{../prompts/eval_qwen.json}
\inputcode[style=logstyle,firstline=241,lastline=320,firstnumber=241,caption={Qwen: итоговый JSON}]{../prompts/eval_qwen.json}
\inputcode[style=logstyle,firstline=321,lastline=400,firstnumber=321,caption={Qwen: итоговый JSON}]{../prompts/eval_qwen.json}
\inputcode[style=logstyle,firstline=401,lastline=480,firstnumber=401,caption={Qwen: итоговый JSON}]{../prompts/eval_qwen.json}
\inputcode[style=logstyle,firstline=481,lastline=560,firstnumber=481,caption={Qwen: итоговый JSON}]{../prompts/eval_qwen.json}
\inputcode[style=logstyle,firstline=561,lastline=617,firstnumber=561,caption={Qwen: итоговый JSON}]{../prompts/eval_qwen.json}
\section{Яндекс.Алиса}
\label{app:json:alice}
\inputcode[style=logstyle,firstline=1,lastline=80,firstnumber=1,caption={Яндекс.Алиса: итоговый JSON}]{../prompts/eval_alice.json}
\inputcode[style=logstyle,firstline=81,lastline=160,firstnumber=81,caption={Яндекс.Алиса: итоговый JSON}]{../prompts/eval_alice.json}
\inputcode[style=logstyle,firstline=161,lastline=240,firstnumber=161,caption={Яндекс.Алиса: итоговый JSON}]{../prompts/eval_alice.json}
\inputcode[style=logstyle,firstline=241,lastline=320,firstnumber=241,caption={Яндекс.Алиса: итоговый JSON}]{../prompts/eval_alice.json}
\inputcode[style=logstyle,firstline=321,lastline=400,firstnumber=321,caption={Яндекс.Алиса: итоговый JSON}]{../prompts/eval_alice.json}
\inputcode[style=logstyle,firstline=401,lastline=480,firstnumber=401,caption={Яндекс.Алиса: итоговый JSON}]{../prompts/eval_alice.json}
\inputcode[style=logstyle,firstline=481,lastline=560,firstnumber=481,caption={Яндекс.Алиса: итоговый JSON}]{../prompts/eval_alice.json}
\inputcode[style=logstyle,firstline=561,lastline=617,firstnumber=561,caption={Яндекс.Алиса: итоговый JSON}]{../prompts/eval_alice.json}
\section{Королёв}
\label{app:json:korolev}
\inputcode[style=logstyle,firstline=1,lastline=80,firstnumber=1,caption={Королёв: итоговый JSON}]{../prompts/eval_korolev.json}
\inputcode[style=logstyle,firstline=81,lastline=160,firstnumber=81,caption={Королёв: итоговый JSON}]{../prompts/eval_korolev.json}
\inputcode[style=logstyle,firstline=161,lastline=240,firstnumber=161,caption={Королёв: итоговый JSON}]{../prompts/eval_korolev.json}
\inputcode[style=logstyle,firstline=241,lastline=320,firstnumber=241,caption={Королёв: итоговый JSON}]{../prompts/eval_korolev.json}
\inputcode[style=logstyle,firstline=321,lastline=400,firstnumber=321,caption={Королёв: итоговый JSON}]{../prompts/eval_korolev.json}
\inputcode[style=logstyle,firstline=401,lastline=480,firstnumber=401,caption={Королёв: итоговый JSON}]{../prompts/eval_korolev.json}
\inputcode[style=logstyle,firstline=481,lastline=560,firstnumber=481,caption={Королёв: итоговый JSON}]{../prompts/eval_korolev.json}
\inputcode[style=logstyle,firstline=561,lastline=617,firstnumber=561,caption={Королёв: итоговый JSON}]{../prompts/eval_korolev.json}
\section{DeepSeek}
\label{app:json:deepseek}
\inputcode[style=logstyle,firstline=1,lastline=80,firstnumber=1,caption={DeepSeek: итоговый JSON}]{../prompts/eval_deepseek.json}
\inputcode[style=logstyle,firstline=81,lastline=160,firstnumber=81,caption={DeepSeek: итоговый JSON}]{../prompts/eval_deepseek.json}
\inputcode[style=logstyle,firstline=161,lastline=240,firstnumber=161,caption={DeepSeek: итоговый JSON}]{../prompts/eval_deepseek.json}
\inputcode[style=logstyle,firstline=241,lastline=320,firstnumber=241,caption={DeepSeek: итоговый JSON}]{../prompts/eval_deepseek.json}
\inputcode[style=logstyle,firstline=321,lastline=400,firstnumber=321,caption={DeepSeek: итоговый JSON}]{../prompts/eval_deepseek.json}
\inputcode[style=logstyle,firstline=401,lastline=480,firstnumber=401,caption={DeepSeek: итоговый JSON}]{../prompts/eval_deepseek.json}
\inputcode[style=logstyle,firstline=481,lastline=560,firstnumber=481,caption={DeepSeek: итоговый JSON}]{../prompts/eval_deepseek.json}
\inputcode[style=logstyle,firstline=561,lastline=617,firstnumber=561,caption={DeepSeek: итоговый JSON}]{../prompts/eval_deepseek.json}
\prilojenie{Итоговая проверка и сводка результатов}
\label{app:summary}
\inputcode[style=logstyle,firstline=1,lastline=27,firstnumber=1,caption={Сводка результатов тестирования}]{../prompts/results.txt}

\end{document}
